\section*{Chapitre \ref{ch:m1:option-contracts-and-exchanges} --- Contrats d'options et bourses d'options}

\begin{solution}{exo:m1:option-contracts-and-exchanges:1}
Le put XYZ expirant le 15 janvier 2027 avec un strike de \$47.50. Le call~:
\texttt{XYZ\textvisiblespace\textvisiblespace\textvisiblespace270319C00112500}.
\end{solution}

\begin{solution}{exo:m1:option-contracts-and-exchanges:2}
À 56~: $10 \times 100 \times (6 - 2.40) = +\$3\,600$. À 48~: les call expirent
sans valeur, $-\$2\,400$. Point mort à $50 + 2.40 = 52.40$.
\end{solution}

\begin{solution}{exo:m1:option-contracts-and-exchanges:3}
$2 \times (11 \times 70 + 5 \times 30) = 1\,840$ séries. Sur seize bourses~:
29\,440 cotations bilatérales à rafraîchir pour un centime sur une action, ce qui explique pourquoi
le market making d'options est d'abord un problème de débit de messages et de choix des
cotations à ne pas actualiser.
\end{solution}

\begin{solution}{exo:m1:option-contracts-and-exchanges:4}
47.00~: les deux put vendus sont assignés, $+5\,000$ actions. 49.00~: seuls les
put 50, $+1\,000$. 50.00~: rien n'est \omterm{def:m1:option-contracts-and-exchanges:series}{dans la monnaie}, 0. 50.01~: les call longs sont
exercés, $+2\,500$. 52.50~: toujours $+2\,500$. 52.51~: les call 60 vendus sont
assignés, $2\,500 - 6\,000 = -3\,500$.
\end{solution}

\begin{solution}{exo:m1:option-contracts-and-exchanges:5}
Oui. Exercer permet de capter le dividende de 0.60 et renonce à 0.15 de valeur temps~;
ne pas exercer conserve une option dont le sous-jacent ouvrira 0.60 plus bas. Le
vendeur couvert par des actions en portefeuille est assigné pendant la nuit~: les actions sont cédées
à 40, le dividende revient au nouveau propriétaire, et le vendeur a perdu le
0.60 sur lequel il comptait, à moins que le prix de l'option n'intègre déjà
l'exercice anticipé. Chaque jour, les desks établissent un rapport sur les exercices avant chaque \omterm{def:m1:shares-corporate-actions-indices:dividend}{date de détachement}.
\end{solution}

\begin{solution}{exo:m1:option-contracts-and-exchanges:6}
Le premier cas~: l'option ne peut plus être négociée après le jeudi, alors que sa valeur est fixée
par les cours d'ouverture du vendredi. Le desk porte l'écart de la nuit et la
dispersion des cours d'ouverture sans pouvoir ajuster l'option, seulement
sa couverture, et doit dénouer la couverture \emph{dans les enchères d'ouverture}. Le
second cas~: l'option se négocie jusqu'à la clôture, donc le risque est continu, mais son
gamma dans les dernières minutes est extrême près du strike, et la couverture doit être
dénouée dans l'enchère de clôture.
\end{solution}

\begin{solution}{exo:m1:option-contracts-and-exchanges:7}
$+2\,500$ actions avec un exercice par exception, 0 avec l'instruction contraire.
Refuser l'exercice est judicieux lorsque le centime de \omterm{def:m1:option-contracts-and-exchanges:series}{valeur intrinsèque} est inférieur au coût
et au risque des actions reçues~: commissions et frais sur 2\,500 actions, et
surtout l'exposition entre la clôture du vendredi et l'ouverture du lundi sur des actions achetées
à 50 qui peuvent ouvrir à 49.50. Pour un détenteur non couvert par une vente à découvert, un centime
ne rémunère pas un week-end.
\end{solution}

\begin{solution}{exo:m1:option-contracts-and-exchanges:8}
Le vendeur a supposé que les porteurs exercent d'après la clôture de 16:00.
Ils disposent jusqu'à l'heure limite du soir et exerceront si l'action se traite
au-dessus de 52.50 hors séance. Le vendeur ne connaît sa position que
le lendemain matin~; avec une couverture dimensionnée pour une option «~expirée sans valeur~», il peut se retrouver vendeur
de 6\,000 actions le lundi. À 15:59, les call auraient pu être rachetés pour
quelques centimes~: la prime de la certitude.
\end{solution}

\begin{solution}{pb:m1:option-contracts-and-exchanges:1}
\textbf{1.} Seuls les call 50 longs, à 2.48. Les call 52.50 sont hors de la
monnaie de 0.02~; les deux séries de put sont hors de la monnaie.
\textbf{2.} Les 25 call longs sont exercés~; tout le reste expire.
\textbf{3.} $+2\,500$ actions, achetées à 50.
\textbf{4.} $-1\,900 + 2\,500 = +600$ actions.
\textbf{5.} La couverture visait un portefeuille qui n'existe plus~: vendre 600 actions,
idéalement dans l'enchère de clôture du vendredi, en estimant à 15:55 ce que l'exercice
livrera.
\textbf{6.} Jusqu'à l'heure limite des bourses, 16:30 à Chicago, 17:30 à New
York, et plus tôt chez la plupart des courtiers.
\textbf{7.} Parce que l'action vaut 53.20 là où elle peut être vendue dès maintenant, et
le call la livre à 52.50~: l'exercice vaut 0.70 par action quel que soit le
cours constaté à 16:00.
\textbf{8.} 4\,500 actions livrées à 52.50.
\textbf{9.} $-1\,900 + 2\,500 - 4\,500 = -3\,900$ actions.
\textbf{10.} $4\,500 \times (54.10 - 52.50) = \$7\,200$.
\textbf{11.} $60 \times 100 \times 0.04 = \$240$.
\textbf{12.} Trente fois moins.
\textbf{13.} Précisément pour cette raison~: l'option n'expire pas à 16:00
mais à l'heure limite, et son vendeur est rémunéré pour quatre-vingt-dix minutes d'informations et pour
le droit de l'acheteur de décider ensuite.
\textbf{14.} Elle couvre le mauvais scénario une fois sur deux~: si l'action
baisse hors séance, personne n'exerce, et le desk se retrouve long de 3\,000 actions indésirables
à une ouverture plus basse. Face aux 54.10 du lundi, elle aurait gagné
\$4\,860 sur les titres tout en restant vendeur de 1\,500 actions. Racheter les options
supprime l'incertitude~; acheter des actions revient à parier sur son issue.
\textbf{15.} Le desk détient ce droit~: il décide, avec le même recul de quatre-vingt-dix
minutes, et peut refuser.
\textbf{16.} Un \omterm{def:m1:what-a-trading-firm-does:market-maker}{teneur de marché} long d'options près du strike est long de gamma~: il
vend des actions lorsque le cours monte au-dessus du strike et en achète lorsqu'il baisse en dessous,
ce qui ramène le cours vers le strike~; l'effet est le plus fort lorsque
les \omterm{def:m1:futures-contracts-and-exchanges:oi}{positions ouvertes} sur ce strike sont importantes par rapport au volume du jour.
\textbf{17.} Non~: rien n'est livré, l'exercice est automatique sur une
valeur de règlement publiée, et il n'y a aucune décision a posteriori. Le
risque se déplace dans la valeur de règlement elle-même (exercice 6).
\textbf{18.} Avec le client, en position vendeuse d'actions acquises au
strike pendant le week-end~; si le client ne peut pas couvrir la marge, avec le
courtier, puis au-delà du courtier avec la garantie de la chambre de compensation.
\textbf{19.} \textbf{Vendeur de 3\,900 actions.}
\textbf{20.} Le porteur de l'option, après la clôture, avec une information que le
vendeur possède aussi mais ne peut pas exploiter.
\end{solution}

\begin{solution}{iq:m1:option-contracts-and-exchanges:1}
Les options américaines peuvent être exercées chaque jour jusqu'à l'\omterm{def:m1:option-contracts-and-exchanges:option}{échéance}, les européennes seulement à
l'\omterm{def:m1:option-contracts-and-exchanges:option}{échéance}. Aux États-Unis, les options sur actions et sur fonds sont américaines et réglées physiquement~;
les principales options sur indices sont européennes et \omterm{def:m1:basis-roll-and-delivery:delivery}{réglées en espèces}. La
différence compte autour des dividendes et pour les puts profondément \omterm{def:m1:option-contracts-and-exchanges:series}{dans la monnaie}, où
l'exercice anticipé est optimal, ainsi que pour le risque d'\omterm{def:m1:option-contracts-and-exchanges:assignment}{assignation} du vendeur.
\iqlookfor{la correspondance des produits et une conséquence.}
\end{solution}

\begin{solution}{iq:m1:option-contracts-and-exchanges:2}
Toutes les options cotées aux États-Unis sont émises et compensées par une seule chambre de compensation, si bien qu'une série
est le même contrat où qu'elle soit négociée~: une position longue sur une bourse et une position courte sur une autre se compensent à zéro. Un contrat à terme
est compensé par la chambre de compensation de sa propre bourse, tandis qu'un contrat identique négocié ailleurs relève d'une autre
contrepartie.
counterparty.
\iqlookfor{la compensation, et non la négociation, source de la fongibilité.}
\end{solution}

\begin{solution}{iq:m1:option-contracts-and-exchanges:3}
Un call~: seulement pour capter un dividende, juste avant la \omterm{def:m1:shares-corporate-actions-indices:dividend}{date de détachement}, lorsque le
dividende dépasse la valeur temps restante (en gros, la valeur du put au même strike
augmentée des intérêts sur le strike). Sinon, vendre le call
vaut mieux que l'exercer. Un put~: lorsqu'il est assez profondément \omterm{def:m1:option-contracts-and-exchanges:series}{dans la monnaie} pour que
les intérêts perçus sur le strike, encaissés aujourd'hui, dépassent la valeur temps restante~;
plus souvent lorsque les taux sont élevés.
\iqlookfor{la comparaison avec la valeur temps dans les deux cas.}
\end{solution}

\begin{solution}{iq:m1:option-contracts-and-exchanges:4}
Être vendeur d'options qui expirent alors que le sous-jacent se trouve sur le strike~: le
nombre d'options assignées, donc la position en actions du lundi, reste inconnu jusqu'au
lendemain matin, et les détenteurs décident après la clôture avec des informations plus récentes. Il faut la gérer
en rachetant avant la clôture les options vendues proches du strike, en neutralisant les strikes qui
expirent durant les derniers jours, en estimant les probabilités d'\omterm{def:m1:option-contracts-and-exchanges:assignment}{assignation} pour ce qui reste, et en
conservant la capacité de négocier après la clôture
et à l'ouverture du lundi.
\iqlookfor{clôturer la position plutôt que couvrir l'inconnue.}
\end{solution}

\begin{solution}{iq:m1:option-contracts-and-exchanges:5}
Un enregistrement immuable~: identifiant du sous-jacent, date d'\omterm{def:m1:option-contracts-and-exchanges:option}{échéance}, droit, strike
sous forme d'entier dans la plus petite unité, multiplicateur, style d'exercice,
type de règlement-livraison, et l'actif livrable pour les contrats ajustés. Jamais de strike flottant~;
jamais le symbole d'affichage comme clé, car les racines changent après les opérations sur le capital.
Il faut associer les symboles de bourse et les chaînes OSI à cet enregistrement au moyen des données de référence,
avec des dates d'effet~; donner à chaque enregistrement un identifiant entier compact pour le
chemin critique.
\iqlookfor{entier pour le strike, contrats ajustés, datation effective.}
\end{solution}

\begin{solution}{iq:m1:option-contracts-and-exchanges:6}
L'option cesse d'être cotée le jeudi~; son paiement dépend d'une cotation construite
à partir du cours d'ouverture de chaque action. Détenir le delta en contrats à terme ou dans un panier
pendant la nuit~; l'exposition qui subsiste est le gamma face au mouvement de nuit,
qui ne peut pas être recouvert dans l'option. Le vendredi matin, la couverture doit
être convertie exactement en valeur de règlement-livraison~: vendre (ou acheter) le panier
aux enchères d'ouverture, ou utiliser des contrats à terme qui se règlent sur la même cotation, ce qui
supprime entièrement la base. Dimensionner la position le jeudi en tenant compte
de l'écart de nuit, et non de la volatilité intrajournalière.
\iqlookfor{faire correspondre la sortie de la couverture au mécanisme de règlement-livraison.}
\end{solution}
